% Original pages: 81--88
% Figures: none
% Uncertainty log: none

iii) Every locally closed subset of $\operatorname{Spec}(A)$ consisting of a single point is closed.

[ii) and iii) are geometrical formulations of conditions ii) and iii) of Exercise 23.]

\noindent\textit{Valuation rings and valuations}

\begin{enumerate}
\setcounter{enumi}{26}
\item Let $A$, $B$ be two local rings. $B$ is said to \textit{dominate} $A$ if $A$ is a subring of $B$ and the maximal ideal $\mathfrak m$ of $A$ is contained in the maximal ideal $\mathfrak n$ of $B$ (or, equivalently, if $\mathfrak m=\mathfrak n\cap A$). Let $K$ be a field and let $\Sigma$ be the set of all local subrings of $K$. If $\Sigma$ is ordered by the relation of domination, show that $\Sigma$ has maximal elements and that $A\in\Sigma$ is maximal if and only if $A$ is a valuation ring of $K$. [Use \amref{5.21}.]

\item Let $A$ be an integral domain, $K$ its field of fractions. Show that the following are equivalent:

(1) $A$ is a valuation ring of $K$;

(2) If $\mathfrak a$, $\mathfrak b$ are any two ideals of $A$, then either $\mathfrak a\subseteq\mathfrak b$ or $\mathfrak b\subseteq\mathfrak a$.

Deduce that if $A$ is a valuation ring and $\mathfrak p$ is a prime ideal of $A$, then $A_{\mathfrak p}$ and $A/\mathfrak p$ are valuation rings of their fields of fractions.

\item Let $A$ be a valuation ring of a field $K$. Show that every subring of $K$ which contains $A$ is a local ring of $A$.

\item Let $A$ be a valuation ring of a field $K$. The group $U$ of units of $A$ is a subgroup of the multiplicative group $K^*$ of $K$.

Let $\Gamma=K^*/U$. If $\xi$, $\eta\in\Gamma$ are represented by $x$, $y\in K$, define $\xi\geq\eta$ to mean $xy^{-1}\in A$. Show that this defines a total ordering on $\Gamma$ which is compatible with the group structure (i.e., $\xi\geq\eta\Rightarrow\xi\omega\geq\eta\omega$ for all $\omega\in\Gamma$). In other words, $\Gamma$ is a totally ordered abelian group. It is called the \textit{value group} of $A$.

Let $v:K^*\longrightarrow\Gamma$ be the canonical homomorphism. Show that $v(x+y)\geq\min(v(x),v(y))$ for all $x,y\in K^*$.

\item Conversely, let $\Gamma$ be a totally ordered abelian group (written additively), and let $K$ be a field. A \textit{valuation of $K$ with values in $\Gamma$} is a mapping $v:K^*\longrightarrow\Gamma$ such that

(1) $v(xy)=v(x)+v(y)$,

(2) $v(x+y)\geq\min(v(x),v(y))$,

for all $x,y\in K^*$. Show that the set of elements $x\in K^*$ such that $v(x)\geq0$ is a valuation ring of $K$. This ring is called the \textit{valuation ring} of $v$, and the subgroup $v(K^*)$ of $\Gamma$ is the \textit{value group} of $v$.

Thus the concepts of valuation ring and valuation are essentially equivalent.

\item Let $\Gamma$ be a totally ordered abelian group. A subgroup $\Delta$ of $\Gamma$ is \textit{isolated} in $\Gamma$ if, whenever $0\leq\beta\leq\alpha$ and $\alpha\in\Delta$, we have $\beta\in\Delta$. Let $A$ be a valuation ring of a field $K$, with value group $\Gamma$ (Exercise 31). If $\mathfrak p$ is a prime ideal of $A$, show that $v(A-\mathfrak p)$ is the set of elements $\geq0$ in an isolated subgroup $\Delta$ of $\Gamma$, and that the mapping so defined of $\operatorname{Spec}(A)$ into the set of isolated subgroups of $\Gamma$ is bijective.

If $\mathfrak p$ is a prime ideal of $A$, what are the value groups of the valuation rings $A/\mathfrak p$, $A_{\mathfrak p}$?

\item Let $\Gamma$ be a totally ordered abelian group. We shall show how to construct a field $K$ and a valuation $v$ of $K$ with $\Gamma$ as value group. Let $k$ be any field and let $A=k[\Gamma]$ be the group algebra of $\Gamma$ over $k$. By definition, $A$ is freely generated as a $k$-vector space by elements $x_\alpha$ ($\alpha\in\Gamma$) such that $x_\alpha x_\beta=x_{\alpha+\beta}$. Show that $A$ is an integral domain.

If $u=\lambda_1x_{\alpha_1}+\cdots+\lambda_nx_{\alpha_n}$ is any non-zero element of $A$, where the $\lambda_i$ are all $\neq0$ and $\alpha_1<\cdots<\alpha_n$, define $v_0(u)$ to be $\alpha_1$. Show that the mapping $v_0:A-\{0\}\longrightarrow\Gamma$ satisfies conditions (1) and (2) of Exercise 31.

Let $K$ be the field of fractions of $A$. Show that $v_0$ can be uniquely extended to a valuation $v$ of $K$, and that the value group of $v$ is precisely $\Gamma$.

\item Let $A$ be a valuation ring and $K$ its field of fractions. Let $f:A\longrightarrow B$ be a ring homomorphism such that $f^*:\operatorname{Spec}(B)\longrightarrow\operatorname{Spec}(A)$ is a closed mapping. Then if $g:B\longrightarrow K$ is any $A$-algebra homomorphism (i.e., if $g\circ f$ is the embedding of $A$ in $K$) we have $g(B)=A$.

[Let $C=g(B)$; obviously $C\supseteq A$. Let $\mathfrak n$ be a maximal ideal of $C$. Since $f^*$ is closed, $\mathfrak m=\mathfrak n\cap A$ is the maximal ideal of $A$, whence $A_{\mathfrak m}=A$. Also the local ring $C_{\mathfrak n}$ dominates $A_{\mathfrak m}$. Hence by Exercise 27 we have $C_{\mathfrak n}=A$ and therefore $C\subseteq A$.]

\item From Exercises 1 and 3 it follows that, if $f:A\longrightarrow B$ is integral and $C$ is any $A$-algebra, then the mapping $(f\otimes1)^*:\operatorname{Spec}(B\otimes_A C)\longrightarrow\operatorname{Spec}(C)$ is a closed map.

Conversely, suppose that $f:A\longrightarrow B$ has this property and that $B$ is an integral domain. Then $f$ is integral. [Replacing $A$ by its image in $B$, reduce to the case where $A\subseteq B$ and $f$ is the injection. Let $K$ be the field of fractions of $B$ and let $A'$ be a valuation ring of $K$ containing $A$. By \amref{5.22} it is enough to show that $A'$ contains $B$. By hypothesis $\operatorname{Spec}(B\otimes_A A')\longrightarrow\operatorname{Spec}(A')$ is a closed map. Apply the result of Exercise 34 to the homomorphism $B\otimes_A A'\longrightarrow K$ defined by $b\otimes a'\mapsto ba'$. It follows that $ba'\in A'$ for all $b\in B$ and all $a'\in A'$; taking $a'=1$, we have what we want.]

Show that the result just proved remains valid if $B$ is a ring with only finitely many minimal prime ideals (e.g., if $B$ is Noetherian). [Let $\mathfrak p_i$ be the minimal prime ideals. Then each composite homomorphism $A\longrightarrow B\longrightarrow B/\mathfrak p_i$ is integral, hence $A\longrightarrow\prod(B/\mathfrak p_i)$ is integral, hence $A\longrightarrow B/\mathfrak N$ is integral (where $\mathfrak N$ is the nilradical of $B$), hence finally $A\longrightarrow B$ is integral.]
\end{enumerate}

\section{Chain Conditions}

So far we have considered quite arbitrary commutative rings (with identity). To go further, however, and obtain deeper theorems we need to impose some finiteness conditions. The most convenient way is in the form of ``chain conditions''. These apply both to rings and modules, and in this chapter we consider the case of modules. Most of the arguments are of a rather formal kind and because of this there is a symmetry between the ascending and descending chains---a symmetry which disappears in the case of rings as we shall see in subsequent chapters.

Let $\Sigma$ be a set partially ordered by a relation $\leq$ (i.e., $\leq$ is reflexive and transitive and is such that $x\leq y$ and $y\leq x$ together imply $x=y$).

\noindent\textit{Proposition 6.1.\amtarget{6.1} The following conditions on $\Sigma$ are equivalent:}
\begin{enumerate}
\renewcommand{\labelenumi}{\roman{enumi})}
\item \textit{Every increasing sequence $x_1\leq x_2\leq\cdots$ in $\Sigma$ is stationary (i.e., there exists $n$ such that $x_n=x_{n+1}=\cdots$).}
\item \textit{Every non-empty subset of $\Sigma$ has a maximal element.}
\end{enumerate}

\noindent\textit{Proof.} i) $\Rightarrow$ ii). If ii) is false there is a non-empty subset $T$ of $\Sigma$ with no maximal element, and we can construct inductively a non-terminating strictly increasing sequence in $T$.

ii) $\Rightarrow$ i). The set $(x_m)_{m\geq1}$ has a maximal element, say $x_n$. \proofbox

If $\Sigma$ is the set of submodules of a module $M$, ordered by the relation $\subseteq$, then i) is called the \textit{ascending chain condition} (a.c.c. for short) and ii) the \textit{maximal condition}. A module $M$ satisfying either of these equivalent conditions is said to be \textit{Noetherian} (after Emmy Noether). If $\Sigma$ is ordered by $\supseteq$, then i) is the \textit{descending chain condition} (d.c.c. for short) and ii) the \textit{minimal condition}. A module $M$ satisfying these is said to be \textit{Artinian} (after Emil Artin).

\noindent\textbf{Examples.} 1) A finite abelian group (as $\mathbb Z$-module) satisfies both a.c.c. and d.c.c.

2) The ring $\mathbb Z$ (as $\mathbb Z$-module) satisfies a.c.c. but not d.c.c. For if $a\in\mathbb Z$ and $a\neq0$ we have $(a)\supset(a^2)\supset\cdots\supset(a^n)\supset\cdots$ (strict inclusions).

3) Let $G$ be the subgroup of $\mathbb Q/\mathbb Z$ consisting of all elements whose order is a power of $p$, where $p$ is a fixed prime. Then $G$ has exactly one subgroup $G_n$ of order $p^n$ for each $n\geq0$, and $G_0\subset G_1\subset\cdots\subset G_n\subset\cdots$ (strict inclusions) so that $G$ does not satisfy the a.c.c. On the other hand the only proper subgroups of $G$ are the $G_n$, so that $G$ does satisfy d.c.c.

4) The group $H$ of all rational numbers of the form $m/p^n$ ($m,n\in\mathbb Z$, $n\geq0$) satisfies neither chain condition. For we have an exact sequence $0\longrightarrow\mathbb Z\longrightarrow H\longrightarrow G\longrightarrow0$, so that $H$ doesn't satisfy d.c.c. because $\mathbb Z$ doesn't; and $H$ doesn't satisfy a.c.c. because $G$ doesn't.

5) The ring $k[x]$ ($k$ a field, $x$ an indeterminate) satisfies a.c.c. but not d.c.c. on ideals.

6) The polynomial ring $k[x_1,x_2,\ldots]$ in an infinite number of indeterminates $x_n$ satisfies neither chain condition on ideals: for the sequence $(x_1)\subset(x_1,x_2)\subset\cdots$ is strictly increasing, and the sequence $(x_1)\supset(x_1^2)\supset(x_1^3)\supset\cdots$ is strictly decreasing.

7) We shall see later that a ring which satisfies d.c.c. on ideals must also satisfy a.c.c. on ideals. (This is \textit{not} true in general for modules: see Examples 2, 3 above.)

\noindent\textit{Proposition 6.2.\amtarget{6.2} $M$ is a Noetherian $A$-module $\Longleftrightarrow$ every submodule of $M$ is finitely generated.}

\noindent\textit{Proof.} $\Rightarrow$: Let $N$ be a submodule of $M$, and let $\Sigma$ be the set of all finitely generated submodules of $N$. Then $\Sigma$ is not empty (since $0\in\Sigma$) and therefore has a maximal element, say $N_0$. If $N_0\neq N$, consider the submodule $N_0+Ax$ where $x\in N$, $x\notin N_0$; this is finitely generated and strictly contains $N_0$, so we have a contradiction. Hence $N=N_0$ and therefore $N$ is finitely generated.

$\Leftarrow$: Let $M_1\subseteq M_2\subseteq\cdots$ be an ascending chain of submodules of $M$. Then $N=\bigcup_{n=1}^{\infty}M_n$ is a submodule of $M$, hence is finitely generated, say by $x_1,\ldots,x_r$. Say $x_i\in M_{n_i}$ and let $n=\max_{i=1}^r n_i$; then each $x_i\in M_n$, hence $M_n=M$ and therefore the chain is stationary. \proofbox

Because of \amref{6.2}, Noetherian modules are more important than Artinian modules: the Noetherian condition is just the right finiteness condition to make a lot of theorems work. However, many of the elementary formal properties apply equally to Noetherian and Artinian modules.

\noindent\textit{Proposition 6.3.\amtarget{6.3} Let $0\longrightarrow M'\xrightarrow{\alpha}M\xrightarrow{\beta}M''\longrightarrow0$ be an exact sequence of $A$-modules. Then}
\begin{enumerate}
\renewcommand{\labelenumi}{\roman{enumi})}
\item \textit{$M$ is Noetherian $\Longleftrightarrow M'$ and $M''$ are Noetherian;}
\item \textit{$M$ is Artinian $\Longleftrightarrow M'$ and $M''$ are Artinian.}
\end{enumerate}

\noindent\textit{Proof.} We shall prove i); the proof of ii) is similar.

$\Rightarrow$: An ascending chain of submodules of $M'$ (or $M''$) gives rise to a chain in $M$, hence is stationary.

$\Leftarrow$: Let $(L_n)_{n\geq1}$ be an ascending chain of submodules of $M$; then $(\alpha^{-1}(L_n))$ is a chain in $M'$, and $(\beta(L_n))$ is a chain in $M''$. For large enough $n$ both these chains are stationary, and it follows that the chain $(L_n)$ is stationary. \proofbox

\noindent\textit{Corollary 6.4.\amtarget{6.4} If $M_i$ $(1\leq i\leq n)$ are Noetherian (resp. Artinian) $A$-modules, so is $\bigoplus_{i=1}^nM_i$.}

\noindent\textit{Proof.} Apply induction and \amref{6.3} to the exact sequence
\[
 0\longrightarrow M_n\longrightarrow\bigoplus_{i=1}^nM_i\longrightarrow\bigoplus_{i=1}^{n-1}M_i\longrightarrow0.
\]
\proofbox

A ring $A$ is said to be Noetherian (resp. Artinian) if it is so as an $A$-module, i.e., if it satisfies a.c.c. (resp. d.c.c.) on ideals.

\noindent\textbf{Examples.} 1) Any field is both Artinian and Noetherian; so is the ring $\mathbb Z/(n)$ ($n\neq0$). The ring $\mathbb Z$ is Noetherian, but not Artinian (Exercise 2 before \amref{6.2}).

2) Any principal ideal domain is Noetherian (by \amref{6.2}: every ideal is finitely generated).

3) The ring $k[x_1,x_2,\ldots]$ is not Noetherian (Exercise 6 above). But it is an integral domain, hence has a field of fractions. Thus a subring of a Noetherian ring need not be Noetherian.

4) Let $X$ be a compact infinite Hausdorff space, $C(X)$ the ring of real-valued continuous functions on $X$. Take a strictly decreasing sequence $F_1\supset F_2\supset\cdots$ of closed sets in $X$, and let $\mathfrak a_n=\{f\in C(X):f(F_n)=0\}$. Then the $\mathfrak a_n$ form a strictly increasing sequence of ideals in $C(X)$: so $C(X)$ is not a Noetherian ring.

\noindent\textit{Proposition 6.5.\amtarget{6.5} Let $A$ be a Noetherian (resp. Artinian) ring, $M$ a finitely-generated $A$-module. Then $M$ is Noetherian (resp. Artinian).}

\noindent\textit{Proof.} $M$ is a quotient of $A^n$ for some $n$: apply \amref{6.4} and \amref{6.3}. \proofbox

\noindent\textit{Proposition 6.6.\amtarget{6.6} Let $A$ be Noetherian (resp. Artinian), $\mathfrak a$ an ideal of $A$. Then $A/\mathfrak a$ is a Noetherian (resp. Artinian) ring.}

\noindent\textit{Proof.} By \amref{6.3} $A/\mathfrak a$ is Noetherian (resp. Artinian) as an $A$-module, hence also as an $A/\mathfrak a$-module. \proofbox

A \textit{chain} of submodules of a module $M$ is a sequence $(M_i)$ ($0\leq i\leq n$) of submodules of $M$ such that
\[
 M=M_0\supset M_1\supset\cdots\supset M_n=0\quad\text{(strict inclusions)}.
\]
The \textit{length} of the chain is $n$ (the number of ``links''). A \textit{composition series} of $M$ is a maximal chain, that is one in which no extra submodules can be inserted: this is equivalent to saying that each quotient $M_{i-1}/M_i$ $(1\leq i\leq n)$ is \textit{simple} (that is, has no submodules except $0$ and itself).

\noindent\textit{Proposition 6.7.\amtarget{6.7} Suppose that $M$ has a composition series of length $n$. Then every composition series of $M$ has length $n$, and every chain in $M$ can be extended to a composition series.}

\noindent\textit{Proof.} Let $l(M)$ denote the least length of a composition series of a module $M$. ($l(M)=+\infty$ if $M$ has no composition series.)

i) $N\subset M\Rightarrow l(N)<l(M)$. Let $(M_i)$ be a composition series of $M$ of minimum length, and consider the submodules $N_i=N\cap M_i$ of $N$. Since $N_{i-1}/N_i\subseteq M_{i-1}/M_i$ and the latter is a simple module, we have either $N_{i-1}/N_i=M_{i-1}/M_i$, or else $N_{i-1}=N_i$; hence, removing repeated terms, we have a composition series of $N$, so that $l(N)\leq l(M)$. If $l(N)=l(M)=n$, then $N_{i-1}/N_i=M_{i-1}/M_i$ for each $i=1,2,\ldots,n$; hence $M_{n-1}=N_{n-1}$, hence $M_{n-2}=N_{n-2}$, \ldots, and finally $M=N$.

ii) \textit{Any chain in $M$ has length $\leq l(M)$.} Let $M=M_0\supset M_1\supset\cdots$ be a chain of length $k$. Then by i) we have $l(M)>l(M_1)>\cdots>l(M_k)=0$, hence $l(M)\geq k$.

iii) Consider any composition series of $M$. If it has length $k$, then $k\leq l(M)$ by ii), hence $k=l(M)$ by the definition of $l(M)$. Hence all composition series have the same length. Finally, consider any chain. If its length is $l(M)$ it must be a composition series, by ii); if its length is $<l(M)$ it is not a composition series, hence not maximal, and therefore new terms can be inserted until the length is $l(M)$. \proofbox

\noindent\textit{Proposition 6.8.\amtarget{6.8} $M$ has a composition series $\Longleftrightarrow M$ satisfies both chain conditions.}

\noindent\textit{Proof.} $\Rightarrow$: All chains in $M$ are of bounded length, hence both a.c.c. and d.c.c. hold.

$\Leftarrow$: Construct a composition series of $M$ as follows. Since $M=M_0$ satisfies the maximum condition by \amref{6.1}, it has a maximal submodule $M_1\subset M_0$. Similarly $M_1$ has a maximal submodule $M_2\subset M_1$, and so on. Thus we have a strictly descending chain $M_0\supset M_1\supset\cdots$ which by d.c.c. must be finite, and hence is a composition series of $M$. \proofbox

A module satisfying both a.c.c. and d.c.c. is therefore called a \textit{module of finite length}. By \amref{6.7} all composition series of $M$ have the same length $l(M)$, called the \textit{length} of $M$. The Jordan--Hölder theorem applies to modules of finite length: if $(M_i)_{0\leq i\leq n}$ and $(M_i')_{0\leq i\leq n}$ are any two composition series of $M$, there is a one-to-one correspondence between the set of quotients $(M_{i-1}/M_i)_{1\leq i\leq n}$ and the set of quotients $(M_{i-1}'/M_i')_{1\leq i\leq n}$, such that corresponding quotients are isomorphic. The proof is the same as for finite groups.

\noindent\textit{Proposition 6.9.\amtarget{6.9} The length $l(M)$ is an additive function on the class of all $A$-modules of finite length.}

\noindent\textit{Proof.} We have to show that if $0\longrightarrow M'\xrightarrow{\alpha}M\xrightarrow{\beta}M''\longrightarrow0$ is an exact sequence, then $l(M)=l(M')+l(M'')$. Take the image under $\alpha$ of any composition series of $M'$ and the inverse image under $\beta$ of any composition series of $M''$; these fit together to give a composition series of $M$, hence the result. \proofbox

Consider the particular case of modules over a field $k$, i.e., $k$-vector spaces:

\noindent\textit{Proposition 6.10.\amtarget{6.10} For $k$-vector spaces $V$ the following conditions are equivalent:}
\begin{enumerate}
\renewcommand{\labelenumi}{\roman{enumi})}
\item \textit{finite dimension;}
\item \textit{finite length;}
\item \textit{a.c.c.;}
\item \textit{d.c.c.}
\end{enumerate}
\textit{Moreover, if these conditions are satisfied, length $={}$ dimension.}

\noindent\textit{Proof.} i) $\Rightarrow$ ii) is elementary; ii) $\Rightarrow$ iii), ii) $\Rightarrow$ iv) from \amref{6.8}. Remains to prove iii) $\Rightarrow$ i) and iv) $\Rightarrow$ i). Suppose i) is false, then there exists an infinite sequence $(x_n)_{n\geq1}$ of linearly independent elements of $V$. Let $U_n$ (resp. $V_n$) be the vector space spanned by $x_1,\ldots,x_n$ (resp., $x_{n+1},x_{n+2},\ldots$). Then the chain $(U_n)_{n\geq1}$ (resp. $(V_n)_{n\geq1}$) is infinite and strictly ascending (resp. strictly descending).

\noindent\textit{Corollary 6.11.\amtarget{6.11} Let $A$ be a ring in which the zero ideal is a product $\mathfrak m_1\cdots\mathfrak m_n$ of (not necessarily distinct) maximal ideals. Then $A$ is Noetherian if and only if $A$ is Artinian.}

\noindent\textit{Proof.} Consider the chain of ideals $A\supseteq\mathfrak m_1\supseteq\mathfrak m_1\mathfrak m_2\supseteq\cdots\supseteq\mathfrak m_1\cdots\mathfrak m_n=0$. Each factor $\mathfrak m_1\cdots\mathfrak m_{i-1}/\mathfrak m_1\cdots\mathfrak m_i$ is a vector space over the field $A/\mathfrak m_i$. Hence a.c.c. $\Longleftrightarrow$ d.c.c. for each factor. But a.c.c. (resp. d.c.c.) for each factor $\Longleftrightarrow$ a.c.c. (resp. d.c.c.) for $A$, by repeated application of \amref{6.3}. Hence a.c.c. $\Longleftrightarrow$ d.c.c. for $A$. \proofbox

\subsection{Exercises}

\begin{enumerate}
\item
\begin{enumerate}
\renewcommand{\labelenumii}{\roman{enumii})}
\item Let $M$ be a Noetherian $A$-module and $u:M\longrightarrow M$ a module homomorphism. If $u$ is surjective, then $u$ is an isomorphism.
\item If $M$ is Artinian and $u$ is injective, then again $u$ is an isomorphism.
\end{enumerate}
[For (i), consider the submodules $\operatorname{Ker}(u^n)$; for (ii), the quotient modules $\operatorname{Coker}(u^n)$.]

\item Let $M$ be an $A$-module. If every non-empty set of finitely generated submodules of $M$ has a maximal element, then $M$ is Noetherian.

\item Let $M$ be an $A$-module and let $N_1$, $N_2$ be submodules of $M$. If $M/N_1$ and $M/N_2$ are Noetherian, so is $M/(N_1\cap N_2)$. Similarly with Artinian in place of Noetherian.

\item Let $M$ be a Noetherian $A$-module and let $\mathfrak a$ be the annihilator of $M$ in $A$. Prove that $A/\mathfrak a$ is a Noetherian ring.

If we replace ``Noetherian'' by ``Artinian'' in this result, is it still true?

\item A topological space $X$ is said to be Noetherian if the open subsets of $X$ satisfy the ascending chain condition (or, equivalently, the maximal condition). Since closed subsets are complements of open subsets, it comes to the same thing to say that the closed subsets of $X$ satisfy the descending chain condition (or, equivalently, the minimal condition). Show that, if $X$ is Noetherian, then every subspace of $X$ is Noetherian, and that $X$ is quasi-compact.

\item Prove that the following are equivalent:
\begin{enumerate}
\renewcommand{\labelenumii}{\roman{enumii})}
\item $X$ is Noetherian.
\item Every open subspace of $X$ is quasi-compact.
\item Every subspace of $X$ is quasi-compact.
\end{enumerate}

\item A Noetherian space is a finite union of irreducible closed subspaces. [Consider the set $\Sigma$ of closed subsets of $X$ which are not finite unions of irreducible closed subspaces.] Hence the set of irreducible components of a Noetherian space is finite.

\item If $A$ is a Noetherian ring then $\operatorname{Spec}(A)$ is a Noetherian topological space. Is the converse true?

\item Deduce from Exercise 8 that the set of minimal prime ideals in a Noetherian ring is finite.

\item If $M$ is a Noetherian module (over an arbitrary ring $A$) then $\operatorname{Supp}(M)$ is a closed Noetherian subspace of $\operatorname{Spec}(A)$.

\item Let $f:A\longrightarrow B$ be a ring homomorphism and suppose that $\operatorname{Spec}(B)$ is a Noetherian space (Exercise 5). Prove that $f^*:\operatorname{Spec}(B)\longrightarrow\operatorname{Spec}(A)$ is a closed mapping if and only if $f$ has the going-up property (Chapter 5, Exercise 10).

\item Let $A$ be a ring such that $\operatorname{Spec}(A)$ is a Noetherian space. Show that the set of prime ideals of $A$ satisfies the ascending chain condition. Is the converse true?
\end{enumerate}
